% --- BEGIN FEATURE ABLATION STUDY ---
\section{Feature ablation study}

\subsection*{Repeat the classification using only feature 1 to 4. What do you obtain?}

When restricting the feature space to the first four features (Features 1, 2, 3, and 4) and discarding the final two, the validation error rates for the three generative classifiers are as follows:
\begin{itemize}
    \item \textbf{Standard MVG:} 8.0\%
    \item \textbf{Naive Bayes:} 7.7\%
    \item \textbf{Tied Covariance MVG:} 9.5\%
\end{itemize}

\subsection*{What can we conclude on discarding the last two features?}

Comparing these results with our full 6-dimensional baseline model, we observe a clear degradation in classification performance across all three architectures. Specifically, the error rate of the unconstrained Standard MVG increases from 7.0\% to 8.0\%, while the Naive Bayes model degrades from 7.2\% to 7.7\%. From this degradation, we can definitively conclude that discarding the last two features results in a net loss of discriminative information, proving that Features 5 and 6 provide valuable class-separating utility when evaluating the full feature vector.

\subsection*{Despite the inaccuracy of the assumption for these two features, are the Gaussian models still able to extract some useful information?}

Yes, the Gaussian models are unequivocally capable of extracting useful, predictive information from Features 5 and 6. As demonstrated in Laboratory 4, the 1D Maximum Likelihood density estimation showed that Features 5 and 6 are heavily skewed to the left and exhibit sharp drop-offs, making them a poor fit for a symmetric Gaussian distribution. 

However, the fact that including these two features improves the full MVG error rate by a absolute 1.0\% (dropping from 8.0\% to 7.0\%) demonstrates a vital machine learning principle: a model can tolerate a significant statistical mismatch if the feature provides strong, independent boundary information. The structural information gained by mapping the data along these two dimensions far outweighs the approximation errors introduced by the Gaussian model's rigid symmetry.

\subsection*{Repeat the classification using only features 1-2 and 3-4. Which model is better in each case?}

Slicing the dataset into isolated 2-dimensional feature pairs reveals a complete reversal of model superiority depending on the selected subspace:
\begin{itemize}
    \item \textbf{Case 1 (Features 1–2 jointly):} The Standard MVG and Naive Bayes models vastly outperform the Tied model. Standard MVG achieves an error rate of \textbf{36.5\%} and Naive Bayes reaches \textbf{36.3\%}, whereas the Tied Covariance MVG model collapses entirely, resulting in an error rate of \textbf{49.4\%}.
    \item \textbf{Case 2 (Features 3–4 jointly):} The performance among the classifiers becomes nearly uniform, but the Tied Covariance MVG emerges as the superior model, achieving the lowest error rate at \textbf{9.4\%}, slightly outperforming both the Standard MVG (9.5\%) and Naive Bayes (9.5\%).
\end{itemize}

\subsection*{How is this related to the characteristics of the two classifiers?}

This compelling behavior is directly tied to the geometric constraints governing the decision boundaries of the respective classifiers. The Standard MVG estimates separate, class-specific covariance matrices, which mathematically translates into a complex, quadratic decision boundary capable of enveloping a class with a larger spread. 

Conversely, the Tied MVG forces a single, pooled within-class covariance matrix across both classes. This constraint restricts its decision boundary to a strictly linear hyperplane. The performance in each subspace depends entirely on whether the physical separation of the classes is driven by their differences in coordinate means or their differences in directional variance.

\subsection*{Is the tied model effective at all for the first two features? Why?}

No, the Tied model is completely ineffective when restricted to the first two features, as its 49.4\% error rate is practically identical to a random coin flip. This total failure is explained by our exploratory analysis from Laboratories 2 and 4. For Features 1 and 2, the empirical means of the "Fake" and "Genuine" classes are located extremely close to one another, centered near zero. The true separation between the classes is driven entirely by their variances: for Feature 1, the Genuine variance ($1.4302$) is more than double the Fake variance ($0.5696$), whereas for Feature 2, the pattern flips and the Fake variance ($1.4209$) dominates the Genuine variance ($0.5783$). 

Because the Tied model operates under the strict assumption that both classes share the exact same variance, it mathematically strips away the class-specific spreads. Since the means are overlapping, forcing equal covariances completely removes the primary source of discriminative information available in that subspace, leaving the linear boundary unable to separate the data.

\subsection*{And the MVG? And for the second pair of features?}

The Standard MVG is highly effective for the first pair of features because its quadratic decision boundary is specifically designed to separate classes based on differences in variance, allowing it to easily draw a boundary that isolates the tightly clustered "Fake" class from the widely dispersed "Genuine" class.

For the second pair of features (Features 3 and 4), the underlying data characteristics are completely reversed. As shown in our MLE logs, the classes share nearly identical variances for these features (e.g., Feature 3 variance is $0.5500$ for Fake and $0.5489$ for Genuine), but their means are strongly separated on opposite sides of the origin (Fake mean is $-0.6809$ and Genuine mean is $+0.6652$). 

Because the true class variances in this subspace are actually equal, the underlying assumption of the Tied MVG model is a perfect reflection of physical reality. A simple linear hyperplane is the mathematically optimal boundary to separate these two parallel, equal-variance Gaussian globs. While the unconstrained Standard MVG is flexible enough to approximate this linear boundary, it wastes parameters estimating two separate covariance matrices from the validation data. By estimating only a single, shared covariance matrix, the Tied MVG eliminates estimation noise and avoids overfitting, which explains why it achieves a superior, highly robust error rate of 9.4\%.
% --- END FEATURE ABLATION STUDY ---